Sigui $A=\begin{pmatrix}1&2\\0&1\end{pmatrix}$.

\begin{apartats}

\apartat[1,25]{0,75}
Comprova que $A$ no commuta amb $B=\begin{pmatrix}1&0\\1&1\end{pmatrix}$.

\begin{solucio}
\[
AB=\begin{pmatrix}3&2\\1&1\end{pmatrix},\qquad BA=\begin{pmatrix}1&2\\1&3\end{pmatrix}.
\]
$AB\ne BA$: no commuten.
\end{solucio}

\begin{tria}{matrius-que-commuten}
\itemtria{totes}{1}{1,25}
Troba totes les matrius quadrades d'ordre 2 que commuten amb $A$.

\begin{solucio}
Sigui $X=\begin{pmatrix}a&b\\c&d\end{pmatrix}$. Aleshores
\[
AX=\begin{pmatrix}a+2c&b+2d\\c&d\end{pmatrix},\qquad XA=\begin{pmatrix}a&2a+b\\c&2c+d\end{pmatrix}.
\]
Igualant: $a+2c=a$ dona $c=0$; $b+2d=2a+b$ dona $d=a$; $d=2c+d$ torna a donar $c=0$. Les matrius que
commuten amb $A$ són les de la forma $X=\begin{pmatrix}a&b\\0&a\end{pmatrix}$, amb $a,b\in\mathbb R$.
\end{solucio}

\itemtria{un-parametre}{1}{1,25}
Troba el valor de $k$ perquè $A$ commuti amb $M=\begin{pmatrix}k&3\\0&2\end{pmatrix}$.

\begin{solucio}
\[
AM=\begin{pmatrix}k&3+4\\0&2\end{pmatrix}=\begin{pmatrix}k&7\\0&2\end{pmatrix},\qquad
MA=\begin{pmatrix}k&2k+3\\0&2\end{pmatrix}.
\]
Només difereixen en l'element de la fila 1 i la columna 2: cal $7=2k+3$, és a dir, $k=2$.
\end{solucio}
\end{tria}

\begin{nomesllarg}
\apartat{0,75}
Demostra que, si una matriu $B$ commuta amb $A$, aleshores $B$ també commuta amb $A+5I$.

\begin{solucio}
$B(A+5I)=BA+5B$ i $(A+5I)B=AB+5B$. Com que $AB=BA$, les dues expressions són iguals. (La identitat
commuta amb qualsevol matriu: $BI=B=IB$.)
\end{solucio}
\end{nomesllarg}

\end{apartats}
